\begin{problem}[Residue field from the quotient domain]
For $\mathfrak p\in\Spec A$, prove
\[
\kappa(\mathfrak p)=A_{\mathfrak p}/\mathfrak pA_{\mathfrak p}
\cong\Frac(A/\mathfrak p).
\]
Use the formula to compute residue fields at $(0)$ and $(t-a)$ in $k[t]$, at $(p)$ in $\ZZ$, and at $(t^2+1)$ in $\RR[t]$.
\end{problem}

\paragraph{Why this problem matters.}
This is the main computational theorem for residue fields and is repeatedly used later in fibers and tangent spaces.

\paragraph{Useful ideas before starting.}
Localize the exact sequence $0\to\mathfrak p\to A\to A/\mathfrak p\to0$.

\paragraph{Guided hints.}
The image of $A\setminus\mathfrak p$ in $A/\mathfrak p$ is the set of nonzero elements.

\begin{solution}
Localization gives
\[
A_{\mathfrak p}/\mathfrak pA_{\mathfrak p}
\cong S^{-1}(A/\mathfrak p),\qquad S=A\setminus\mathfrak p.
\]
Since $A/\mathfrak p$ is a domain and $S$ maps onto its nonzero elements, the right side is its fraction field.  Therefore
\[
\kappa((0)\subset k[t])=k(t),\quad
\kappa((t-a))=k,
\]
\[
\kappa((p)\subset\ZZ)=\FF_p,
\quad
\kappa((t^2+1)\subset\RR[t])=\CC.
\]
\end{solution}

\paragraph{What happened mathematically.}
Closed rational points, closed non-rational points, and generic points are unified by one formula.

\paragraph{Extensions.}
Compute residue fields of all closed points of $\Spec\FF_q[t]$.

\paragraph{Provenance.}
Promotes the central residue-field theorem and its worked examples from legacy \emph{Theory of Commutative Algebra 12} into a full mastery dossier; the source theorem was not originally numbered as a Problem.
